\section{AI Teammate Interaction and Impact}
\label{sec:ai-report}

\sectionnote{Owner: Derek Goodman; every member supplies AI-use logs. Target: exactly the last 2 pages. Rubric: Interaction with AI teammate; Quality and impact of AI teammate.}

\begin{requirements}
  \item ``Great'' on the interaction criterion means AI model selection, task assignment and the AI interaction protocol, including the prompting strategy, are explained well enough that another group could reproduce them. ``Great'' on the impact criterion means a clear and reasonable quality-control procedure, an insightful reflection on human--AI team dynamics, and transparency about the AI's quantitative impact on the report.
  \item Build it from the interaction logs every member keeps from the first prompt, and submit the logs with the report.
  \item Show that AI use followed a deliberate plan rather than being ad hoc.
\end{requirements}

\subsection{AI Member Profile and Selection Process}
\label{sec:ai-profile}

\begin{requirements}
  \item The name and version of every AI tool used.
  \item Why each was the right choice for A2's tasks, compared with at least one alternative the group considered. In A1 we wrote that we did not compare tools; the rubric asks for model selection to be explained.
\end{requirements}

\subsection{Tasks Assigned to the AI Teammate}
\label{sec:ai-tasks}

\begin{requirements}
  \item An itemised list of the tasks given to the AI, who assigned each, and why it suited an AI. Likely A2 tasks: summarising the git history behind a divergence, explaining Rust traits and modules found in Understand, proposing mapping rules for files that are hard to place, and drafting or revising text.
\end{requirements}

\subsection{Interaction Protocol and Prompting Strategy}
\label{sec:ai-protocol}

\begin{requirements}
  \item Who prompted and the workflow, described well enough that another group could reproduce it.
  \item At least one complex or critical prompt, quoted, and how it was refined (context given, persona, output format).
\end{requirements}

\subsection{Validation and Quality Control Procedures}
\label{sec:ai-validation}

\begin{requirements}
  \item The specific checks applied to AI output, for example: every dependency the AI claimed was confirmed in Understand; every commit hash, author and date was confirmed in a clone at \pinnedcommit; every method name in a sequence diagram was confirmed in the source.
  \item Who carried out each check, and one example of an AI error the checks caught.
\end{requirements}

\subsection{Quantitative Contribution to the Final Deliverable}
\label{sec:ai-contribution}

\begin{requirements}
  \item A percentage estimate of the AI's contribution to the submitted report, with a table breaking it down by area, and how the estimate was made (for example, from the logs).
  \item The estimate must hold up against the submitted logs; the rubric marks down an impact that is ``missing or highly inaccurate''.
\end{requirements}

\subsection{Reflection on Human--AI Team Dynamics}
\label{sec:ai-reflection}

\begin{requirements}
  \item The AI's strengths and limitations on this deliverable: whether it saved time or created new work, how it influenced decisions, and any disagreements it caused.
  \item What changed in how we worked with the AI since A1, and what we will keep for A3.
\end{requirements}
